\documentclass[11pt]{article}
\usepackage[a4paper,margin=2.05cm]{geometry}
\usepackage[T1]{fontenc}
\usepackage{lmodern}
\usepackage{amsmath,amssymb}
\usepackage{xcolor}
\usepackage{microtype}
\usepackage[hidelinks]{hyperref}
\setlength{\emergencystretch}{3em}

\definecolor{ink}{HTML}{172033}
\definecolor{accent}{HTML}{385B7A}
\definecolor{soft}{HTML}{EEF5FA}
\color{ink}

\title{A complete characterization of discretizable continuous frames}
\author{Literature answer to the discretization problem in arXiv:math/0410571}
\date{11 August 2026}

\begin{document}
\maketitle

\begin{center}
\fcolorbox{accent}{soft}{\parbox{0.91\textwidth}{
\textbf{Status: literature-already-answered, full characterization.}
A measurable family can be sampled to a discrete Hilbert frame exactly when it is an essentially bounded continuous frame for some positive $\sigma$-finite measure.  Hence every bounded continuous frame is sampleable, but arbitrary unbounded continuous frames need not be.}}
\end{center}

\section*{Source problem}
A continuous frame $\Psi:X\to H$ with respect to a measure $\mu$ satisfies
\[
 A\|h\|^2\le \int_X |\langle h,\Psi(t)\rangle|^2\,d\mu(t)
 \le B\|h\|^2 \qquad(h\in H).
\]
The Ali--Antoine--Gazeau discretization problem asks for conditions under which one may choose a countable sampling sequence $(t_j)$ so that $(\Psi(t_j))$ is a discrete frame.  Fornasier--Rauhut describe their contribution as a partial answer on source PDF pages 1--2 and give sufficient kernel-localization conditions; on source PDF page 31 they explicitly ask whether their $L^1$-integrability hypothesis is necessary \cite{FR}.

\section*{Explicit later solution}
Freeman--Speegle say in their abstract and introduction that they solve the discretization problem in full generality.  Their Theorem 1.3, on supporting PDF pages 2--3, states the following \cite{FS}.

Let $(X,\Sigma)$ have measurable singletons and let $\Psi:X\to H$ be measurable.  The following are equivalent:
\begin{enumerate}
\item There is a countable sequence $(t_j)$ in $X$ (repetitions allowed) such that $(\Psi(t_j))$ is a frame for $H$.
\item There is a positive $\sigma$-finite measure $\nu$ on $(X,\Sigma)$ such that $\Psi$ is a continuous frame with respect to $\nu$ and $\|\Psi(t)\|$ is bounded outside a $\nu$-null set.
\end{enumerate}
In particular every bounded continuous frame is sampleable.  Theorem 1.4 strengthens this uniformly: there are universal $C,D>0$ such that every continuous Parseval frame taking values in the Hilbert unit ball has a sampling with lower and upper frame bounds $C,D$.

The universal statement without boundedness is false.  Freeman--Speegle give the continuous frame
\[
 \Psi(n)=\sqrt n\,e_n\quad(n\in\mathbb N),
 \qquad \mu(\{n\})=1/n,
\]
for $\ell_2$.  Any sampling with dense span must contain every direction $e_n$, and its vector norms are then unbounded; a discrete frame is necessarily norm bounded.  Thus this frame cannot be discretized.

\section*{Proof intuition}
Normalize a bounded continuous frame by its frame operator to a Parseval frame in the unit ball.  The Marcus--Spielman--Srivastava/Kadison--Singer partition machinery, assembled through finite-dimensional approximations, selects countably many frame vectors with universal two-sided frame bounds.  Undoing the normalization yields a sampled frame for the original family.  Conversely, if a sampling already exists, put counting mass (including multiplicity) on the sampled points.  The same family is then an essentially bounded continuous frame for that atomic $\sigma$-finite measure.  These two directions produce the exact criterion.

\section*{Provenance and scope}
This is an explicit later solution, not an agent-generated theorem.  It fully resolves when the underlying Hilbert continuous frame admits a sampled Hilbert frame.  It does not assert that every sampling also yields the source's simultaneous atomic decompositions and Banach frames for all associated coorbit spaces; those stronger conclusions retain localization hypotheses.  Nor does it settle the separate source remark asking whether the kernel algebra $\mathcal A_1$ is spectral.

{\raggedright
\begin{thebibliography}{9}
\bibitem{FR}
M. Fornasier and H. Rauhut,
\emph{Continuous Frames, Function Spaces, and the Discretization Problem},
arXiv:math/0410571v1 (2004).

\bibitem{FS}
D. Freeman and D. Speegle,
\emph{The discretization problem for continuous frames},
arXiv:1611.06469 (2016), Theorems 1.3 and 1.4.
\end{thebibliography}
\par}

\end{document}
